\documentclass[10pt,a4paper]{amsart}
\usepackage[margin=19mm]{geometry}
\usepackage{amsmath,amssymb,mathtools}
\usepackage[scr=rsfs,cal=euler]{mathalfa}
\usepackage{xfrac}
\usepackage[unicode,hidelinks,bookmarks=false]{hyperref}
\newcommand{\ie}{i.e.\ }
\newcommand{\cf}{cf.\ }
\newcommand{\resp}{respectively\ }
\newcommand{\Subsection}{\subsection}
\providecommand{\nobreakdash}{\nobreak\hbox{-}}
\setlength{\emergencystretch}{2em}
\multlinegap0pt
\usepackage{tikz}
\usepackage{amssymb}
\usepackage{enumerate}
\usepackage{float}
\newtheorem{proposition}{Proposition}
\DeclareMathOperator{\N}{N}
\title{Bicyclic graphs with the smallest and largest numbers of connected sets}
\author{Audace A. V. Dossou-Olory}
\date{Source: arXiv:2603.26812v1 (CC BY 4.0)}
\begin{document}
\maketitle

\noindent\emph{Source and licence.} Bicyclic graphs with the smallest and largest numbers of connected sets. Version arXiv:2603.26812v1 (CC BY 4.0), distributed by arXiv under the Creative Commons Attribution 4.0 International licence, http://creativecommons.org/licenses/by/4.0/, as recorded in the arXiv licence metadata for this exact version and shown on the abstract page https://arxiv.org/abs/2603.26812v1. Full licence notice: \texttt{licenses/arxiv-2603.26812-CC-BY-4.0.txt}.\par\smallskip

\noindent\textit{Editorial scope.} Counted unit: the article's proposition Prop:nplus3 (source0.tex lines 331--336). The article's complete original proof of it is delivered with it (source0.tex lines 338--366, one \texttt{proof} environment of 29 lines). No mathematics is written, repaired or re-derived: the statement and the proof are the article's own lines, in the article's own notation, and no retained line is substituted or re-flowed.  No further source-local context is needed: every cross-reference of every retained line resolves inside the two delivered spans.  Nothing was omitted from any retained span, so the package carries no omission record.  No retained line contains a citation, so the wrapper carries no bibliography and no bibliography entry.  No retained line contains a figure environment, an image-inclusion command or drawing code, so no image is delivered and the package declares no asset dependency.  Editorial formatting only, in addition: an AMS \texttt{amsart} preamble that restates the article's own declarations and macros used by the retained lines, namely the environments \texttt{proposition} on separate counters, together with the standard packages those lines need.  The retained environments print in this document as Proposition 1 in order of appearance, whereas the article prints the same items with the numbering of its own full document.  The complete native source of the article as distributed in the arXiv e-print for this exact version is bundled unchanged as \texttt{sources/197304/source0.tex}.
\begin{proposition}\label{Prop:nplus3}
If $G\in \mathbb{B}_n$ and $v\in V(G)$, then 
\begin{align*}
\N(G)_v \geq n+3\,.
\end{align*}
\end{proposition}

\begin{proof}
It is clear that $\N(A_4)_v=7$ if $v$ is of degree two, and  $\N(A_4)_v=8$ if $v$ is of degree three. So the proposition holds (with equality) for $n=4$. Henceforth, assume $n\geq 5$. We distinguish between three scenarios:
\begin{itemize}
\item $v$ is a pendant vertex of $G$, and $v'$ is the neighbour of $v$.\\
Since $G-v$ is also a bicyclic graph, we have
\begin{align*}
\N(G)_v=1+\N(G-v)_{v'} \geq 1+ (n-1+3)=n+3\,,
\end{align*}
where the inequality follows by induction on $n$.
\item $G$ contains a pendant vertex $u \neq v$.\\
Since $G-u$ is also a bicyclic graph, we have
\begin{align*}
\N(G)_v=\N(G-u)_v + \N(G)_{u,v} \geq (n-1+3)  + \N(G)_{u,v} \,,
\end{align*}
where the inequality follows by induction on $n$. The count for $\N(G)_{u,v}$ includes a shortest path between $u$ and $v$, thus $\N(G)_{u,v} \geq 1$. In particular,
\begin{align*}
\N(G)_v \geq  (n-1+3)  + 1 = n+3\,.
\end{align*}
\item $G$ does not have a pendant vertex.\\
All shortest paths starting at $v$ in $G$ contribute to the count $\N(G)_v$. Thus
$$\N(G)_v \geq n\,.$$
Vertex $v$ together with any two of its neighbours also induces a connected graph in which $v$ is not pendant. We deduce the following:
\begin{itemize}
\item If $v$ has at least three neighbours in $G$, then $\N(G)_v \geq n +3 +1$, where the final $1$ accounts for $G$ itself.
\item If $v$ has only two neighbours, say $v_1,v_2$ in $G$, then either $v_1$ or $v_2$ must have a neighbour, say $w$ not belonging to $\{v,v_1,v_2\}$. Thus, each of the graphs induced by $\{v,v_1,v_2\}$ and $\{v,v_1,v_2,w\}$ is connected and none of them has $v$ as a pendant vertex. Hence  $\N(G)_v \geq n +2 +1$, where the final $1$ accounts for the whole $G$.
\end{itemize}
\end{itemize}
This completes the proof of the proposition.
\end{proof}

\end{document}
